\documentclass[../../main.tex]{subfiles}
\begin{document}

\begin{center}
  {\fontsize{48pt}{58pt}\selectfont \textbf{1986}\normalsize\textbf{年度}}
\end{center}

\begin{table}[htb]
  \centering
  \caption{出題テーマと難易度}
  \begin{tabular}{|c|p{9cm}|c|}
    \hline
    \textbf{問題番号} & \textbf{テーマ} & \textbf{難易度} \\
    \hline
    第1問 & 双曲線の下の三角形の最大面積 & \\
    \hline
    第2問 & 平行移動した楕円の共通部分の面積 & \\
    \hline
    第3問 & 平面の法線ベクトルと回転 & \\
    \hline
    第4問 & 二次方程式の係数と実根の範囲 & \\
    \hline
    第5問 & じゃんけんによる移動と勝敗の確率 & \\
    \hline
    第6問 & 傾けた円錐と水面の楕円 & \\
    \hline
  \end{tabular}
\end{table}

\begin{table}[htb]
  \centering
  \caption{解答}
  \renewcommand{\arraystretch}{1.4}
  \begin{tabular}{|c|c|p{9cm}|}
    \hline
    \textbf{問題番号} & \textbf{小問} & \textbf{答え} \\
    \hline
    第1問 & -- & 最大値は$1$．そのとき$(a,b)=(0,2/c)$または$(2c,0)$ \\
    \hline
    第2問 & -- & $M=\dfrac{5\pi}{3}-1$ \\
    \hline
    第3問 & (1) & $\vec{n}_0=\pm\dfrac{1}{3}(1,2,-2)$ \\
    \cline{2-3}
    & (2) & 最小値は$0$，最大値は$\dfrac{1+\sqrt{17}}{6}$ \\
    \hline
    第4問 & (1) & 点$P$の領域は解答の図を参照 \\
    \cline{2-3}
    & (2) & 最小値は$\dfrac{27+3\sqrt{15}}{11}$，最大値は$\dfrac{11+2\sqrt{19}}{3}$ \\
    \hline
    第5問 & (1) & $q=\dfrac{2}{3^n}{}_n\mathrm{C}_k\left(2^{k-1}-1+\dfrac{k}{n}\right)$ \\
    \cline{2-3}
    & (2) & $p(x,y)$の場合分けは解答を参照 \\
    \cline{2-3}
    & (3) & $p=\dfrac{n^4-2n^3+7n^2+2n+4}{4\cdot3^n}$ \\
    \hline
    第6問 & (1) & $\ell=\dfrac{m}{1+m}$，$a=\dfrac{\sqrt{1+m^2}}{1+m}$，$b=\dfrac{\sqrt{1-m^2}}{1+m}$ \\
    \cline{2-3}
    & (2) & $m=\dfrac{2^{2/3}-1}{2^{2/3}+1}$ \\
    \hline
  \end{tabular}
\end{table}

\paragraph{試験形態}
6問．試験時間は要確認．

\end{document}
